\documentclass[12pt,a4paper]{article}
\usepackage[margin=25mm, headheight=15pt, headsep=10pt]{geometry}
\usepackage{fontspec}
\usepackage[table]{xcolor} % Note: table option is required for rowcolors
\usepackage{titlesec}
\usepackage{tocloft}
\usepackage{fancyhdr}
\usepackage{booktabs}
\usepackage{tcolorbox}
\tcbuselibrary{skins, breakable}
\usepackage[colorlinks=true, linkcolor=emerald, urlcolor=emerald, bookmarks=true]{hyperref}

\definecolor{emerald}{RGB}{0,102,51}
\definecolor{darkred}{RGB}{139,0,0}
\definecolor{lightgray}{RGB}{245,245,245}

\usepackage[nil]{babel}
\babelprovide[import, main]{bengali}
\babelprovide[import]{arabic}
\babelfont{rm}[Renderer=HarfBuzz,Scale=1.0]{Noto Serif Bengali}
\babelfont{sf}[Renderer=HarfBuzz]{Noto Sans Bengali}
\babelfont{tt}[Renderer=HarfBuzz]{Noto Sans Bengali}
\babelfont[arabic]{rm}[Renderer=HarfBuzz,Scale=1.1]{Noto Naskh Arabic}

% Section styling
\titleformat{\section}{\Large\bfseries\color{emerald}}{\thesection.}{0.5em}{}
\titleformat{\subsection}{\large\bfseries\color{emerald!85!black}}{\thesubsection.}{0.5em}{}
\titleformat{\subsubsection}{\normalsize\bfseries\color{emerald!70!black}}{\thesubsubsection.}{0.5em}{}
\titlespacing*{\section}{0pt}{18pt}{8pt}

% Fancy Header/Footer setup
\pagestyle{fancy}
\fancyhf{} % clear all header and footer fields
\fancyhead[L]{\color{emerald}\small\textbf{\nouppercase{\leftmark}}}
\fancyfoot[L]{\scriptsize\color{black!50}{\lat{AI}}-সংকলিত, আলেম-যাচাইকৃত নয়}
\fancyfoot[C]{\color{emerald!70!black}\thepage}
\renewcommand{\headrulewidth}{0.4pt}
\renewcommand{\headrule}{\hbox to\headwidth{\color{emerald}\leaders\hrule height \headrulewidth\hfill}}
\renewcommand{\footrulewidth}{0pt}

% Custom boxes
% 1. Ayat/Hadith Box
\newtcolorbox{ayatbox}[1][]{
    enhanced, breakable, 
    colback=emerald!5, 
    colframe=emerald, 
    boxrule=0.5pt, 
    arc=3pt, 
    left=10pt, right=10pt, top=10pt, bottom=10pt,
    #1
}

% 2. Quote/Translation Box
\newtcolorbox{quotebox}[1][]{
    enhanced, breakable,
    frame hidden,
    colback=lightgray,
    borderline west={3pt}{0pt}{emerald},
    left=15pt, right=10pt, top=10pt, bottom=10pt,
    sharp corners,
    #1
}

% 3. AI-generated notice box
\newtcolorbox{warnbox}[1][]{
    enhanced, breakable,
    colback=red!4,
    colframe=darkred,
    boxrule=0.8pt,
    arc=3pt,
    left=10pt, right=10pt, top=10pt, bottom=10pt,
    #1
}

% Commands
\newcommand{\arab}[1]{\foreignlanguage{arabic}{#1}}
\newcommand{\kib}[1]{\textcolor{emerald}{\textbf{#1}}}

% Latin (English) fallback: Noto Serif Bengali has NO Latin glyphs
\newfontfamily\latfont[Renderer=HarfBuzz]{Noto Serif}
\newcommand{\lat}[1]{{\latfont #1}}

\setlength{\parskip}{5pt}
\setlength{\parindent}{0pt}

% Fix for missing bullet in Noto Serif Bengali
\renewcommand{\labelitemi}{$\bullet$}



\begin{document}
\begin{center}{\Large\bfseries\color{emerald} আল্লাহর রহমত ও তাওবা — \lat{07}-বারবার-পাপ-বারবার-ক্ষমা-প্রার্থনা}\end{center}
\vspace{1em}
\section*{ঘটনা ৭ — বারবার পাপ, বারবার ক্ষমা: \lat{“}আমার বান্দা জানে তার রব আছে…\lat{”}}

\begin{warnbox}
 \textbf{গুরুত্বপূর্ণ নোটিশ (\lat{Important} \lat{Notice}):} এই কন্টেন্টটি \lat{AI} (কৃত্রিম বুদ্ধিমত্তা) দিয়ে প্রস্তুত ও সংকলিত। \lat{AI} কঠোর সূত্র-মান নীতি মেনে শুধু নির্ভরযোগ্য উৎস থেকে তথ্য সংগ্রহ করেছে — কেবল সহীহ/হাসান হাদিস ও সহীহ সনদের আসার রাখা হয়েছে, দুর্বল সনদ বাদ দেওয়া হয়েছে। তবে এটি কোনো আলেম বা মুহাদ্দিস কর্তৃক যাচাইকৃত নয়।
\end{warnbox}

\begin{quote}
\textbf{সম্পর্কিত ফাইল:} হাদিস ফাইল — মুসলিম ২৭৫৯ (বারবার পাপ-তাওবা); ব্যবহারিক গাইড — আবার পাপ হলে আবার তাওবা।
\end{quote}

\medskip\hrule\medskip

\subsection*{১. সূত্র কার্ড (\lat{Source} \lat{Card})}

\begin{table}[h]\centering\small
\begin{tabular}{ll}
বিষয় & বিবরণ \\
\hline
\textbf{ঘটনা} & এক বান্দা বারবার পাপ করেছে, বারবার \lat{“}রব, আমাকে ক্ষমা করো\lat{”} বলেছে — আল্লাহ প্রতিবার ক্ষমা করেছেন — কারণ বান্দা জানে তার এমন রব আছে, যিনি ক্ষমা করেন \\
\textbf{বর্ণনাকারী} & আবু হুরাইরা (রাঃ) \\
\textbf{গ্রন্থ ও নম্বর} & সহীহ বুখারি ৭৫০৭; সহীহ মুসলিম ২৭৫৮ \\
\textbf{অধ্যায়} & কিতাবুত তাওহীদ — আল্লাহর ক্ষমা সম্পর্কে \\
\textbf{সনদের মান} & \textbf{সহীহ} (বুখারি-মুসলিম) \\
\textbf{স্থান ও সময়} & মদিনা — নবীজি (সা.) বর্ণনা করেছেন \\
\hline
\end{tabular}
\end{table}

\medskip\hrule\medskip

\subsection*{২. সংক্ষিপ্ত পটভূমি (\lat{Background})}

এটি বনি ইসরাঈলের (অথবা সাধারণ) এক বান্দার ঘটনা — যাকে আল্লাহ ক্ষমা করার পদ্ধতি বর্ণনা করেছেন। হাদিসটির বিশেষত্ব — পাপ ও ক্ষমা প্রার্থনার \textbf{পুনরাবৃত্তি}: পাপ $\rightarrow$ ক্ষমা চাওয়া $\rightarrow$ ক্ষমা; আবার পাপ $\rightarrow$ আবার ক্ষমা চাওয়া $\rightarrow$ আবার ক্ষমা — তিনবার। অর্থ — আল্লাহ পাপের পুনরাবৃত্তিতে তাওবার দরজা বন্ধ করেন না; বরং বারবার ক্ষমা চাওয়ার অভ্যাসকেই তিনি স্বাগত জানান।

\medskip\hrule\medskip

\subsection*{৩. পূর্ণ ঘটনার বিশুদ্ধ বর্ণনা (\lat{The} \lat{Full} \lat{Narration})}

\begin{quote}
নবীজি (সা.) বলেছেন: \textbf{\lat{“}এক বান্দা পাপ করল — সে বলল: হে আমার রব! আমি পাপ করেছি — আমাকে ক্ষমা করো। তার রব বললেন: আমার বান্দা জেনেছে — তার এমন রব আছে, যিনি পাপ ক্ষমা করেন এবং (চাইলে) পাপের জন্য শাস্তিও দেন। আমি আমার বান্দাকে ক্ষমা করে দিলাম। এরপর সে আল্লাহ যতদিন চাইলেন ততদিন (পাপমুক্ত) থাকল, তারপর আবার পাপ করল — বলল: হে আমার রব! আমি পাপ করেছি — আমাকে ক্ষমা করো। তার রব বললেন: আমার বান্দা জেনেছে — তার এমন রব আছে, যিনি পাপ ক্ষমা করেন এবং শাস্তিও দেন। আমি আমার বান্দাকে ক্ষমা করে দিলাম। এরপর সে আবার পাপ করল — বলল: হে আমার রব! আমি পাপ করেছি — আমাকে ক্ষমা করো। তার রব বললেন: আমার বান্দা জেনেছে — তার এমন রব আছে, যিনি পাপ ক্ষমা করেন এবং পাপের জন্য শাস্তিও দেন। আমি আমার বান্দাকে ক্ষমা করে দিলাম — সে যা চায় করুক (এখন তার জন্য ক্ষমা প্রস্তুত)।\lat{”}}
\end{quote}

\medskip\hrule\medskip

\subsection*{৪. শব্দের অর্থ (\lat{Key} \lat{Words})}

\begin{table}[h]\centering\small
\begin{tabular}{ll}
শব্দ & অর্থ \\
\hline
\textbf{আসাবা যাম্বান (\arab{أصاب} \arab{ذنبا})} & পাপ করল — পাপে পড়ল \\
\textbf{আ'লিমা আবদি} & আমার বান্দা জেনেছে — জ্ঞানের প্রশংসা \\
\textbf{ইয়াগফিরুয যাম্বা ওয়া ইয়া'খুজু বিহি} & পাপ ক্ষমা করেন এবং পাপের শাস্তিও দেন — দুই-ই তাঁর ক্ষমতায় \\
\textbf{ফাগাফারতু লিআবদি} & আমি আমার বান্দাকে ক্ষমা করে দিলাম \\
\hline
\end{tabular}
\end{table}

\medskip\hrule\medskip

\subsection*{৫. সংশ্লিষ্ট দলিল ও রেফারেন্স}

\begin{quote}
\textbf{\lat{“}আল্লাহ রাতে তাঁর হাত বাড়ান… যতক্ষণ না পশ্চিম দিক থেকে সূর্য উদিত হয়।\lat{”}} — সহীহ মুসলিম ২৭৫৯ (অধ্যায়: পাপ ও তাওবা বারবার হলেও কবুল)
\end{quote}

\begin{quote}
\textbf{\lat{“}তোমরা রাতে ও দিনে পাপ করতে থাকো, আর আমি সব পাপ ক্ষমা করি — সুতরাং ক্ষমা চাও।\lat{”}} — সহীহ মুসলিম ২৫৭৭
\end{quote}

\begin{quote}
\textbf{\lat{“}হে আদম সন্তান! যতক্ষণ তুমি আমাকে ডাকবে ও আমার কাছে আশা রাখবে, আমি ক্ষমা করব।\lat{”}} — তিরমিযী ৩৫৪০ (হাসান)
\end{quote}

\subsubsection*{মূল শিক্ষা (দলিলভিত্তিক)}

\begin{enumerate}
\item \textbf{ক্ষমার জ্ঞানই ক্ষমার পথ:} বান্দার ক্ষমা পাওয়ার কারণ — সে জানে তার রব ক্ষমা করেন; আল্লাহর ক্ষমার প্রতি এই জ্ঞান ও আশা নিজেই দরজা খোলার চাবি।
\item \textbf{বারবার পাপ $\rightarrow$ বারবার ক্ষমা:} তিনবারের বর্ণনা — আল্লাহ পাপের পুনরাবৃত্তিতে তাওবা বন্ধ করেন না; শর্ত — প্রতিবারই আন্তরিক ক্ষমা প্রার্থনা।
\item \textbf{সতর্কতা:} হাদিসের শেষাংশ \lat{“}সে যা চায় করুক\lat{”} — পাপে উৎসাহ নয়; এটি ক্ষমার ব্যাপকতার বর্ণনা। আলেমরা বলেছেন — এতে তাওবার পরও পাপ করলে নিরাশ না হওয়ার শিক্ষা; পাপে ডুবে থাকার অনুমতি নয় (ব্যবহারিক গাইড ৩ ও ৬ নং)।
\end{enumerate}
\medskip\hrule\medskip

\subsection*{৬. রেফারেন্স}

\begin{itemize}
\item সহীহ বুখারি ৭৫০৭ (কিতাবুত তাওহীদ)
\item সহীহ মুসলিম ২৭৫৮ (কিতাবুত তাওবাহ)
\item কুরআন: নিসা ৪:১১০, যুমার ৩৯:৫৩
\item ব্যাখ্যা: ফাতহুল বারি — এই হাদিসের ব্যাখ্যায় পাপ-পুনরাবৃত্তি ও ক্ষমার আলোচনা
\end{itemize}

\end{document}
